
\documentclass[addpoints]{exam}
\usepackage{amsmath}
\usepackage{amssymb}
\usepackage{color}
\usepackage[version=4]{mhchem}
\usepackage{siunitx}

% \printanswers

\newcommand{\event}{Solutions \& Solubility}
\newcommand{\tournament}{}
\def\type{0}

\setlength\fillinlinelength{\textwidth}
\CorrectChoiceEmphasis{\color{red}\bfseries}
\SolutionEmphasis{\color{red}\bfseries}

\setlength\answerclearance{2pt}
\pagestyle{headandfoot}
\runningheadrule
\runningheader{\event}{\tournament}{\ifnum\type=1 Class Set Test \else Full Name:\hspace{1cm} \fi}
\runningfootrule
\runningfooter{}{Page \thepage\ of \numpages}{}

\begin{document}
\begin{center}
    {\huge Grade 11 Chemistry \\ \event
    \ifprintanswers \space Key
    \else \space \ifnum\type<2 Test \fi \ifnum\type=1 \textbf{(CLASS SET)} \fi
          \ifnum\type=2 Answer Sheet \fi
    \fi \\ \vspace{3mm}\tournament\vspace{1mm}}
\end{center}

\vspace{.25\textheight}

\begin{center}
    First Name: \ifprintanswers
    \underline{\hspace{3.5cm}}\textbf{KEY}\underline{\hspace{5.5cm}}\\\vspace{5mm}
    \else \underline{\hspace{10cm}}\\ \vspace{5mm} \fi
    Last Name: \underline{\hspace{10cm}}\\ \vspace{5mm}
\end{center}

\noindent\textbf{\underline{Directions:}}
\begin{itemize}
    \ifnum\type=1 \item \textbf{This test is a class set.} Do not write on this paper. \fi
    \item Show all work with units. Use $N_A = 6.022 \times 10^{23}$ mol$^{-1}$.
    \item The test is designed to be completed in 75 minutes.
\end{itemize}
\begin{center}
\textit{For grading use only}\\
\multirowgradetable{1}[pages]
\end{center}

\newpage
\begin{questions}

\section*{Part A --- Multiple Choice (15 marks)}
\vspace{5mm}

\question[1] In a solution, the substance present in the largest amount is called the:
\begin{choices}
    \choice Solute
    \correctchoice Solvent
    \choice Electrolyte
    \choice Precipitate
\end{choices}

\question[1] The molar concentration of a solution is defined as:
\begin{choices}
    \choice grams of solute per litre of solution
    \correctchoice moles of solute per litre of solution
    \choice moles of solute per kilogram of solvent
    \choice grams of solute per 100 mL of solution
\end{choices}

\question[1] What is the concentration of a solution prepared by dissolving $\SI{4.0}{mol}$
of \ce{NaCl} in enough water to make $\SI{2.0}{L}$ of solution?
\begin{choices}
    \choice \SI{0.50}{mol/L}
    \correctchoice \SI{2.0}{mol/L}
    \choice \SI{8.0}{mol/L}
    \choice \SI{4.0}{mol/L}
\end{choices}

\question[1] How many moles of solute are in $\SI{250}{mL}$ of a $\SI{0.400}{mol/L}$
solution?
\begin{choices}
    \choice $0.400$ mol
    \choice $0.160$ mol
    \correctchoice $0.100$ mol
    \choice $1.60$ mol
\end{choices}

\question[1] A saturated solution is one that:
\begin{choices}
    \choice Contains the maximum possible amount of solvent
    \correctchoice Contains the maximum amount of dissolved solute at a given temperature
    \choice Has a concentration of exactly \SI{1.0}{mol/L}
    \choice Contains no undissolved solute
\end{choices}

\question[1] Which of the following \textbf{increases} the solubility of a gas in water?
\begin{choices}
    \choice Increasing temperature
    \correctchoice Increasing pressure
    \choice Decreasing pressure
    \choice Stirring
\end{choices}

\question[1] The dilution formula $C_1V_1 = C_2V_2$ is valid because:
\begin{choices}
    \choice Volume is conserved during dilution
    \correctchoice The number of moles of solute is conserved during dilution
    \choice Concentration doubles when volume halves
    \choice Mass is always conserved
\end{choices}

\question[1] A $\SI{500}{mL}$ solution of \SI{2.0}{mol/L} \ce{HCl} is diluted to
\SI{1000}{mL}. The new concentration is:
\begin{choices}
    \choice \SI{4.0}{mol/L}
    \correctchoice \SI{1.0}{mol/L}
    \choice \SI{0.50}{mol/L}
    \choice \SI{2.0}{mol/L}
\end{choices}

\question[1] Which of the following compounds is a strong electrolyte when dissolved in water?
\begin{choices}
    \choice \ce{C6H12O6} (glucose)
    \choice \ce{CH3COOH} (acetic acid)
    \correctchoice \ce{KNO3}
    \choice \ce{C2H5OH} (ethanol)
\end{choices}

\question[1] The process by which water molecules surround and separate ions from an ionic
crystal is called:
\begin{choices}
    \choice Precipitation
    \correctchoice Hydration (solvation)
    \choice Ionization
    \choice Electrolysis
\end{choices}

\question[1] The solubility of most solid solutes in water:
\begin{choices}
    \choice Decreases with increasing temperature
    \correctchoice Increases with increasing temperature
    \choice Is unaffected by temperature
    \choice Depends only on pressure
\end{choices}

\question[1] Which ion concentration is highest in a solution of $\SI{0.10}{mol/L}$
\ce{MgCl2}?
\begin{choices}
    \choice $[\ce{Mg^{2+}}] = \SI{0.20}{mol/L}$
    \correctchoice $[\ce{Cl-}] = \SI{0.20}{mol/L}$
    \choice $[\ce{Mg^{2+}}] = \SI{0.10}{mol/L}$
    \choice $[\ce{Cl-}] = \SI{0.10}{mol/L}$
\end{choices}

\question[1] To prepare $\SI{100.0}{mL}$ of $\SI{0.500}{mol/L}$ \ce{NaOH} from a
$\SI{2.00}{mol/L}$ stock solution, you need:
\begin{choices}
    \choice $\SI{10.0}{mL}$ of stock
    \correctchoice $\SI{25.0}{mL}$ of stock
    \choice $\SI{40.0}{mL}$ of stock
    \choice $\SI{50.0}{mL}$ of stock
\end{choices}

\question[1] An aqueous solution of \ce{NaCl} conducts electricity because:
\begin{choices}
    \choice NaCl molecules carry charge
    \correctchoice \ce{Na+} and \ce{Cl-} ions are free to move
    \choice Water molecules are polar
    \choice Electrons flow through the solution
\end{choices}

\question[1] The number of moles of \ce{Na+} ions in $\SI{150}{mL}$ of
$\SI{0.60}{mol/L}$ \ce{Na2SO4} is:
\begin{choices}
    \choice $0.090$ mol
    \correctchoice $0.18$ mol
    \choice $0.36$ mol
    \choice $0.060$ mol
\end{choices}


\section*{Part B --- Short Answer (33 marks)}

\question Solution preparation and concentration.
\begin{parts}
    \part[2] Calculate the concentration (in mol/L) of a solution made by dissolving
    $\SI{8.55}{g}$ of \ce{KNO3} ($M = \SI{101}{g/mol}$) in enough water to make
    $\SI{500.0}{mL}$ of solution.
    \part[2] What mass of \ce{NaOH} ($M = \SI{40.0}{g/mol}$) is needed to prepare
    $\SI{250}{mL}$ of a $\SI{0.800}{mol/L}$ solution?
    \part[3] What volume of $\SI{6.0}{mol/L}$ \ce{HCl} must be diluted to prepare
    $\SI{300.0}{mL}$ of $\SI{0.50}{mol/L}$ \ce{HCl}? Describe the procedure.
\end{parts}
\begin{solution}[9cm]
\textbf{(a)} $n = \dfrac{8.55}{101} = 0.0847$ mol;\quad
$C = \dfrac{0.0847}{0.5000} = \mathbf{\SI{0.169}{mol/L}}$

\textbf{(b)} $n = C \times V = 0.800 \times 0.250 = 0.200$ mol;\quad
$m = 0.200 \times 40.0 = \mathbf{\SI{8.00}{g}}$

\textbf{(c)} $C_1V_1 = C_2V_2$:\quad
$V_1 = \dfrac{C_2V_2}{C_1} = \dfrac{0.50 \times 0.3000}{6.0} = \mathbf{\SI{25.0}{mL}}$

Procedure: Measure \SI{25.0}{mL} of the \SI{6.0}{mol/L} \ce{HCl} stock. Add it to a
\SI{300}{mL} volumetric flask containing some water. Dilute to the \SI{300}{mL} mark
with distilled water and mix thoroughly.
\end{solution}

\question Ion concentrations.
\begin{parts}
    \part[3] A solution is prepared by dissolving $\SI{14.8}{g}$ of \ce{Ca(NO3)2}
    ($M = \SI{164}{g/mol}$) in enough water to make $\SI{400.0}{mL}$ of solution.
    Calculate the concentration of \ce{Ca^{2+}} ions and \ce{NO3^-} ions.
    \part[2] Two solutions are mixed: $\SI{100.0}{mL}$ of $\SI{0.200}{mol/L}$
    \ce{NaCl} and $\SI{150.0}{mL}$ of $\SI{0.300}{mol/L}$ \ce{NaCl}.
    What is the concentration of \ce{NaCl} in the resulting mixture?
\end{parts}
\begin{solution}[9cm]
\textbf{(a)} $n(\ce{Ca(NO3)2}) = \dfrac{14.8}{164} = 0.09024$ mol in \SI{0.4000}{L}

$[\ce{Ca(NO3)2}] = \dfrac{0.09024}{0.4000} = \SI{0.2256}{mol/L}$

\ce{Ca(NO3)2 -> Ca^{2+} + 2NO3^-}

$[\ce{Ca^{2+}}] = \mathbf{\SI{0.226}{mol/L}}$;\quad
$[\ce{NO3^-}] = 2 \times 0.226 = \mathbf{\SI{0.451}{mol/L}}$

\textbf{(b)} Moles of NaCl: $n_1 = 0.100 \times 0.200 = 0.0200$ mol;\;
$n_2 = 0.150 \times 0.300 = 0.0450$ mol

Total: $n = 0.0650$ mol in $V = \SI{250.0}{mL} = \SI{0.2500}{L}$

$C = \dfrac{0.0650}{0.2500} = \mathbf{\SI{0.260}{mol/L}}$
\end{solution}

\question Solution stoichiometry. Excess \ce{BaCl2(aq)} is added to $\SI{50.0}{mL}$
of $\SI{0.200}{mol/L}$ \ce{Na2SO4(aq)}.
\[\ce{BaCl2(aq) + Na2SO4(aq) -> BaSO4(s) + 2NaCl(aq)}\]
\begin{parts}
    \part[2] Calculate the moles of \ce{Na2SO4} present.
    \part[3] Calculate the mass of \ce{BaSO4} precipitate formed. ($M(\ce{BaSO4}) = \SI{233}{g/mol}$)
    \part[2] Write the net ionic equation for this reaction.
\end{parts}
\begin{solution}[9cm]
\textbf{(a)} $n(\ce{Na2SO4}) = 0.200 \times 0.0500 = \mathbf{0.0100 \text{ mol}}$

\textbf{(b)} From balanced equation, 1 mol \ce{Na2SO4} produces 1 mol \ce{BaSO4}:

$n(\ce{BaSO4}) = 0.0100$ mol;\quad $m = 0.0100 \times 233 = \mathbf{\SI{2.33}{g}}$

\textbf{(c)} Net ionic equation:
\[\ce{Ba^{2+}(aq) + SO4^{2-}(aq) -> BaSO4(s)}\]
\end{solution}

\question Solubility and factors affecting it.
\begin{parts}
    \part[3] A saturated solution of \ce{KNO3} is prepared at $\SI{60}{\celsius}$ and then
    cooled to $\SI{20}{\celsius}$. At $\SI{60}{\celsius}$, the solubility of \ce{KNO3} is
    $\SI{110}{g}$ per $\SI{100}{g}$ water; at $\SI{20}{\celsius}$, it is $\SI{32}{g}$ per
    $\SI{100}{g}$ water. How many grams of \ce{KNO3} crystallize out of $\SI{200}{g}$ of
    water upon cooling?
    \part[2] Explain why carbonated beverages go flat when left open.
    \part[2] Explain why fish are stressed when water temperature rises.
\end{parts}
\begin{solution}[8cm]
\textbf{(a)} In \SI{200}{g} water:
At \SI{60}{\celsius}: $110 \times 2 = \SI{220}{g}$ dissolved.
At \SI{20}{\celsius}: $32 \times 2 = \SI{64}{g}$ remains dissolved.
Mass crystallized: $220 - 64 = \mathbf{\SI{156}{g}}$

\textbf{(b)} \ce{CO2} gas is dissolved under high pressure during carbonation. When the cap
is removed, the pressure above the liquid drops to atmospheric pressure. By Henry's Law,
lower pressure means lower gas solubility, so \ce{CO2} escapes from solution.

\textbf{(c)} The solubility of gases (like \ce{O2}) \textbf{decreases} with increasing
temperature. Warmer water contains less dissolved oxygen, meaning fish cannot extract
enough \ce{O2} through their gills to meet their metabolic needs.
\end{solution}

\end{questions}
\end{document}
